Compiler optimizations are easy to state informally and surprisingly hard to
get exactly right. This is particularly true for \llvm, whose intermediate
representation carries subtle semantic concepts such as \poison,
\vundef, integer overflow flags, and target-independent bitvector behavior.
Production optimization passes such as \instcombine are large, heavily
engineered, and known to be bug-prone~\cite{yang2011finding}. Tools such as
\alive~\cite{lopes2021alive2} demonstrate that formal reasoning can be
practical for many optimization patterns, while projects such as
\vellvm~\cite{zhao2012formalizing} and \leanmlir~\cite{bhat2024verifying}
show that proof assistants can support much deeper semantic developments.

This project explores that space from a deliberately smaller angle. The goal
was not to compete with full \llvm formalizations, but to learn how to build
an executable semantics and machine-checked refinement proofs in \lean for a
carefully chosen fragment. The final artifact is
\proj, a compact \lean library in which one can write
small \llvm-like programs directly in \lean, instantiate symbolic widths,
execute the resulting definitions, and prove refinement theorems between
definitions.

The project is best introduced through the kind of theorems it supports. One
example is a width-generic algebraic identity: for every width~\texttt{w},
\texttt{add i\$0 \%x, 0} is refined by \texttt{return \%x},
where \texttt{i\$0} is a symbolic width and \texttt{\%x} is a symbolic variable.
Another example is more semantic and interesting: for every width~\texttt{w},
\texttt{add i\$0 undef, undef} is refined by \texttt{mul i\$0 undef, 2} in the
project's explicit-\vundef model, meaning every value produced by the
right-hand side is either returned exactly by the left-hand side or justified
by the refinement relation's poison-aware value clause. These statements are not just comments or
unit tests; they are theorems checked by \lean against the executable
semantics.

The final system differs in several ways from the original proposal. The
proposal aimed at a somewhat broader \llvm fragment, including memory
operations such as \texttt{load} and \texttt{store}, and suggested building a
small optimizer function and then proving it sound. Under the available time,
We instead focused on constructing a coherent core: abstract syntax, embedded
surface syntax, width instantiation, well-formedness, semantics,
\vundef handling, and refinement proofs for representative examples. This
decision reduced scope, but it also produced a more complete and better
understood semantic artifact than a partially finished optimizer would have.

The rest of this report explains what was built, how the proofs were carried
out, where the design was intentionally simplified, what parts of the
proposal were not completed, and what we learned from the process.
